
\documentstyle[amsfonts,12pt]{article}

\newtheorem{theorem}{Theorem} 
\newtheorem{corollary}[theorem]{Corollary} 
\newtheorem{conjecture}[theorem]{Conjecture} 
\newtheorem{lemma}[theorem]{Lemma} 
\newtheorem{proposition}[theorem]{Proposition} 
\newtheorem{definition}[theorem]{Definition} 
\newtheorem{example}[theorem]{Example} 
\newtheorem{exercise}[theorem]{Exercise} 
\newtheorem{axiom}{Axiom} 
\newtheorem{remark}{Remark} 
\def\ppp{{\mathbb{P}}}
\def\aaa{{\mathbb{A}}}
\def\fff{{\mathbb{F}}}
\def\qqq{\mathbb{Q}}
\def\rrr{\mathbb{R}}
\def\ccc{\mathbb{C}}
\def\zzz{\mathbb{Z}}
\def\nnn{\mathbb{N}}
\def\pf{{\bf proof}:\ }
\def\qed{$\Box$}

\begin{document}

\author{David Joyner}
\title{Which $S_m$ have a complement in $S_n$?\thanks{Written in 1997, 
slightly revised 2007. wdj@usna.edu.}}
%\date{1-3-2007}
\maketitle

\begin{abstract}
In this paper,
we classify which symmetric groups $S_m$ have a complement in $S_n$, 
$m<n$. Such questions are of interest to computer scientists
who wish for efficient methods to store permutations. 

\end{abstract}

%tableofcontents

This paper is inspired by a question of Herbert Kociemba,
a computer scientist who has one of the best Rubik's
cube solving programs known \cite{K}. Efficient methods of
storing permutations in $S_E$ and $S_V$
(the groups of all permutations of the edges $E$ and 
vertices $V$, respectively, of the Rubik's cube)
are needed, hence leading naturally to the concept
of the complement of $S_m$ in $S_n$. Specifically,
he asked if $S_8$ has a complement in $S_{12}$
(this terminology is defined below). The answer is,
as we shall see, ``no.'' Nonetheless, it turns out to be 
possible to introduce a slightly more general notion
of a ``$k$-tuple of complementary subgroups'' (defined below)
for which the answer to the analogous question is
``yes.''

Notation: If $X$ is any finite set then 
\begin{itemize}
\item
$|X|$ denotes the number of elements in $X$.
\item
$S_X$ denotes the symmetric group on $X$.
\item
$S_n$ denotes the symmetric group on $\{1,2,...,n\}$.
\item
$A_n$ denotes its alternating subgroup of even permutations.
\item
$C_n$ denotes the cyclic subgroup of $S_n$ generated by
the $n$-cycle $(1,2,...,n)$.
\item
 $M_{10}$ denotes ``the Mathieu
group of degree $10$'' and order $720=10!/7!$, which we define
as the subgroup of $S_{10}$ generated by 
$(2,6,10,7)(3,9,4,5)$ and $(1,5)(3,8)(4,10)(7,9)$. 
\item
$M_{11}$ denotes the Mathieu group of degree $11$ and order 
$7920=11!/7!$ generated by $(1,2,3,4,5,6,7,8,9,10,11)$ 
and $(3,7,11,8)(4,10,5,6)$. 
\item
$M_{12}$ denotes the Mathieu group of degree 
$12$ and order 
$95040=12!/7!$ generated by $(2,3,5,9,8,10,6,11,4,7,12)$ and 
\newline
$(1,12)(2,11)(3,10)(4,9)(5,8)(6,7)$. 
\item
For any prime power $q$, ${\Bbb F}_q$ denotes the finite field with
$q$ elements. 
\item
$AGL_n({{\Bbb F}}_q)$ denotes the affine
group of transformations on ${\Bbb F}_q^n$ of the
form $\vec{v}\longmapsto A\vec{v}+\vec{a}$,
where $A\in GL_n({{\Bbb F}}_q)$ and $\vec{a} \in {\Bbb F}_q^n$.
\end{itemize}

For other undefined groups arising below, we refer to \cite{DM}.

\section{Background}

We recall some basic definitions for convenience of the reader.

\subsection{Complements}

If $G$ is a finite group and $G_1,G_2$ are subgroups then we say 
$G_2$ is the {\bf complement} of $G_1$ when
\begin{itemize}
\item
$G_1\cap G_2=\{1\}$, the identity of $G$,
and
\item 
$G=G_1\cdot G_2=\{g_1g_2\ |\ g_1\in G_1,\ g_2\in G_2\}$.
\end{itemize}

If $G_1$ and $G_2$ are complementary subgroups then
we may write each $g\in G$ in the form $g=g_1g_2$,
for some unique $g_1\in G_1$ and some unique $g_2\in G_2$. 
We call $g_1$ and $g_2$ the {\bf factors} of $g$. 
This implies the following result, whose (easy) proof
is omitted.

\begin{lemma}
\label{lemma:projections}
Suppose $G$ is a finite group and $G_2$ is the complement of $G_1$
in the sense above.
\begin{itemize}
\item
There are well-defined projection maps
$p_1:G\rightarrow G_1$, $p_2:G\rightarrow G_2$,
defined by $p_1(g)=g_1$ and $p_2(g)=g_2$, where $g=g_1g_2$
is the factorization of $g$.
\item
$G=G_2\cdot G_1$ and we may write each $g\in G$ in the form $g=g_2g_1$,
for some unique $g_1\in G_1$ and some unique $g_2\in G_2$. In other
words, $G_1$ is the complement of $G_2$ in the sense above.
\end{itemize}
\end{lemma}

As is well-known, for any finite group, the right 
multiplication map $\rho(x):g\longmapsto gx^{-1}$,
$g,x\in G$, yields a homomorphism
$\rho:G\rightarrow S_G$,
where $S_G$ is the symmetric group on $G$. Though it shall not be needed,
it is remarkable that a similar result holds for each of the complements:

\begin{lemma} 
\label{lemma:multiplication_mapping}
Let $G=G_1\cdot G_2$ be a finite group with
$G_2$ a complement of $G_1$. 

(a) The right 
multiplication map $\rho(x):g\longmapsto gx^{-1}$,
$g,x\in G$, yields a homomorphism
\[
\rho_2:G\rightarrow S_{G_2},
\]
where $S_{G_2}$ is the symmetric group on $G_2$,
$\rho_2(x)(g)=p_2(\rho(x)(g))$. 

(b) The left multiplication map $\lambda(x):g\longmapsto xg$,
$g,x\in G$, yields a homomorphism
\[
\lambda_1:G\rightarrow S_{G_1},
\]
$\lambda_1(x)(g)=p_1(\lambda(x)(g))$.
\end{lemma}

\begin{remark} The Lemma above was inspired by a remark of
Kosiemba \cite{K}, in connection with an algorithm for his Rubik's 
cube solving program.
\end{remark}

{\bf proof}: 
Write $g\in G$ in the form $g=g_1g_2$,
for some unique $g_1\in G_1$ and some unique $g_2\in G_2$. 
Fix $x\in G$ and
let $g_1'\in G_1$ and $g_2'\in G_2$ denote the factors
of $gx$: $gx=g_1'g_2'$. 

Claim: $g_2'$ depends on $g_2$ and $x$ but not $g_1$. Indeed,
if $g_2x=g_1''g_2''$, for some unique $g_1''\in G_1$ and 
some unique $g_2''\in G_2$
(which depend only on $x,g_2$ but not on $g_1$),
then $gx=g_1(g_2x)=(g_1g_1'')g_2''$. By uniqueness, $g_2'=g_2''$
and so the claim follows.

Define $r_{x^{-1}}:G_2\rightarrow G_2$ by $r_{x^{-1}}(g_2)=g_2'$.
This is simply the restriction of the right translation map
$\rho(x):G\rightarrow G$ followed by projection onto the $G_2$ factor. 

Claim: $r_{x^{-1}}$ is bijective. Surjectivity is obvious.
For injectivity, if $g_2,h_2\in G_2$
and $g_2x=g_1'g_2'$, $h_2x=g_1''g_2'$ are factorizations with
the same $G_2$ factor then plugging $g_2'=(g_1')^{-1}g_2x$
into the second equation gives $h_2x=(g_1''(g_1')^{-1})g_2x$,
which implies $h_2g_2^{-1}=g_1''(g_1')^{-1}$. This is
impossible (by uniqueness) unless $h_2=g_2$, which proves the claim.

Thus we have a map
\[
\rho_2 : G\rightarrow S_{G_2}
\]
defined by $\rho_2(x)=r_{x^{-1}}$.

Claim: $\rho_2$ is a homomorphism. Indeed, if we fix $x,y\in G$ and
$g_1\in G_1$, $g_2\in G_2$ then we have
\[
\begin{array}{c}
g_1'''r_{(xy)^{-1}}(g_2)=g_1'r_{x^{-1}}(g_2)y
=g_1''r_{y^{-1}}(r_{x^{-1}}(g_2)),
\end{array}
\]
for some $g_1',g_1'',g_1'''\in G_1$. By uniqueness, we must have
\[
r_{(xy)^{-1}}=r_{y^{-1}}\circ r_{x^{-1}},
\]
which implies $\rho_2$ is a homomorphism. 

This completes the proof of the first statement 
(for right multiplication) of the Lemma.
The proof of the second statement (for left multiplication)
is similar and hence omitted.

$\Box$

We end this discussion of complements with the following
generalization:

\begin{definition} 
\label{def:k-compl}
Let $G_1,G_2,...,G_k$ denote subgroups
of $G$. More generally, we say that these subgroups are
{\bf complementary} if
\begin{itemize}
\item
for all $1\leq i<j\leq k$, we have
$G_i\cap G_j=\{1\}$, the identity of $G$,
\item
for all $1\leq i< k$, we have
$G_i\cap G_{i+1}\cdot ...\cdot G_k=\{1\}$, the identity of $G$,
\item
for all $1\leq i< k$, $G_{i+1}\cdot ...\cdot G_k$ is a group,
and
\item 
$G=G_1\cdot ...\cdot G_k=\{g_1g_2...g_k\ |\ g_i\in G_i,\ 1\leq i\leq k\}$.
\end{itemize}
\end{definition}
 
\begin{remark}
Lemma \ref{lemma:projections} above has an 
``obvious'' generalization to the case of 
$k$ complementary subgroups. Lemma 
\ref{lemma:multiplication_mapping} (b) also has an analog.
The statement and proof of these
are left to the reader.

We also point out that 
\[
|G|=|G_1|\cdot ... \cdot |G_k|.
\]
One can use this to show that if $I\subset \{1,2,...,k\}$ is any subset
and if we define the {\it subset} $G_I\subset G$ by
\[
G_I=\prod_{i\in I}G_i=\{g_{i_1}\cdot ...\cdot g_{i_r}
\ |\ {\rm where}\ g_i\in G_i,\ i\in I=\{i_1,...,i_r\}\},
\]
then we have
\[
G_I\cap G_{I^c}=\{1\},
\]
where $I^c=\{1,2,...,k\}-I$. 
\end{remark}

\subsection{Multiply transitive actions}

Let $X$ denote a finite set. If $G$ is a subgroup of $S_X$
and $x\in X$ then we let $G_x$ denote the {\bf stabilizer} of
$x$ in $G$:
\[
G_x=\{ g\in G\ |\ g(x)=x\}.
\]

\begin{definition} 
Let $G$ be a permutation group acting on a 
finite set $X$ (so $G$ is a subgroup of the symmetric 
group of $X$, $S_X$). Let $k\geq 1$ be an integer and let
\[
X^{(k)}=\{{\rm distinct\ }k{\rm -tuples\ in\ }X\}
=\{(x_1,x_2,...,x_k)\ |\ x_i\not= x_j,\ 1\leq i<j\leq k\}.
\]
We say $G$ acts {\bf $k$-transitively on $X$} if $G$ acts
transitively on $X^{(k)}$ via the ``diagonal'' action
$g:(x_1,x_2,...,x_k)\longmapsto (g(x_1),g(x_2),...,g(x_k))$.

If $G$ acts transitively on $X$ and $G_x=1$ for some
(hence all) $x\in X$ then we say $G$ acts {\bf regularly}
on $X$. If $G$ acts $k$-transitively on $X$ 
and acts regularly on $X^{(k)}$ then we say $G$ acts
{\bf sharply $k$-transitively} on $X$.
\end{definition}

Cameron \cite{C}, Dixon and Morimer \cite{DM}, Conway,
Hulpke, and McKay \cite{CHM} are good references for
transitive groups.

\begin{example} If $G$ is a subgroup of $S_n$
and $x=(n-k+1,...,n)\in X^{(k)}$ then $G_x \cong S_{n-k}$.
\end{example}

Though it will not be needed here, we give the well-known
classification of $k$-transitive groups ($k\geq 4$)
due to Jordan below.

\begin{theorem} (Jordan)
 A sharply k-transitive group, $k\geq 4$, must be one of the 
following.

\begin{itemize}
\item
$k \geq  6$: $S_k$ , $S_{k+1}$ and $A_{k+2}$ only.
\item
$k = 5$: $S_5$ , $S_6$ , $A_7$ and the Mathieu group $M_{12}$.
\item
$k = 4$: $S_4$ , $S_5$ , $A_6$ and the Mathieu group $M_{11}$.
\end{itemize}
\end{theorem}

The classification of multiply transitive groups was completed
(for $k = 2$ and $3$) by Zassenhaus in the 1930's. 
For the precise statements for $k = 2$ and $3$, see
\cite{C} or \cite{DM} (or the proof of Theorem \ref{mainthrm} below).

\section{A table}

We give a table which indicates, for small values of $n$, which
$S_m$ have a complement in $S_n$. 

\vskip .6in

\begin{center}
\begin{tabular}{|c|c|c|} \hline
$n$ & $m$ & complement of $S_m$ in $S_n$ \\ \hline
$4$ & $2$ & $A_4$\\ \hline
$4$ & $3$ & $C_4$\\ \hline
$5$ & $2$ & $A_5$\\ \hline
$5$ & $3$ & $\langle (1,2,3,4,5),(2,3,5,4)\rangle 
\cong AGL_1({{\Bbb F}}_5)$ \\
 & & size $20$\\ \hline
$5$ & $4$ & $C_5$\\ \hline

$6$ & $2$ & $A_6$\\ \hline
$6$ & $3$ & $\langle (2,3,4,5,6),(3,4,6,5)\rangle 
\cong PGL_2({{\Bbb F}}_5)$ \\
 & & size $120$\\ \hline
$6$ & $5$ & $C_6$\\ \hline
\end{tabular}
\end{center}

\newpage

\begin{center}
\begin{tabular}{|c|c|c|} \hline
$7$ & $2$ & $A_7$\\ \hline
$7$ & $5$ & $\langle (1,2,6,5,3,7),(1,4,3,2,5,6)\rangle 
\cong AGL_1({{\Bbb F}}_7)$ \\
 & & size $42$\\ \hline
$7$ & $6$ & $C_7$\\ \hline
$8$ & $2$ & $A_8$\\ \hline

$8$ & $5$ & $\langle (1,5,8,6,7,4),(1,5,8,7)(2,3,4,6)\rangle 
\cong PGL_2({{\Bbb F}}_7)$ \\
 & & size $336$\\ \hline
$8$ & $6$ & $AGL_1({{\Bbb F}}_8)$ \\
 & & size $56$\\ \hline
$8$ & $7$ & $C_8$\\ \hline
$9$ & $2$ & $A_9$\\ \hline
$9$ & $6$ & $PGL_2({{\Bbb F}}_8)$ \\
 & & size $504$\\ \hline
$9$ & $7$ & $AGL_1({{\Bbb F}}_9)$ \\
 & & size $72$\\ \hline
$9$ & $8$ & $C_9$\\ \hline
$10$ & $2$ & $A_{10}$\\ \hline
$10$ & $7$ & $M_{10}$ \\
 & & size $720$\\ 
 & & ($PGL_2({{\Bbb F}}_9)$ is another) \\ \hline
$10$ & $9$ & $C_{10}$\\ \hline

$n$ & $m$ & complement of $S_m$ in $S_n$ \\ \hline

$11$ & $2$ & $A_{11}$\\ \hline
$11$ & $7$ & $M_{11}$ \\
 & & size $7920$\\ \hline
$11$ & $9$ & $AGL_1({{\Bbb F}}_{11})$ \\
 & & size $110$\\ \hline
$11$ & $10$ & $C_{11}$\\ \hline
$12$ & $2$ & $A_{12}$\\ \hline
$12$ & $7$ & $M_{12}$ \\
 & & size $95040$\\ \hline
$12$ & $9$ & $ PGL_2({{\Bbb F}}_{11})$ \\
 & & $=\langle (1,2,3,4,5,6,7,8,9,10,11),
(1,2,4,8,5,10,9,7,3,6),$\\
 & & $(1,10)(2,5)(3,7)(4,8)(6,9)(11,12)\rangle $ \\
 & & size $1320$\\ \hline
$12$ & $11$ & $C_{12}$\\ \hline
\end{tabular}
\end{center}

\section{General results}

The following result is a special case of Jordan's theorem
above (see also, for example, \cite{R}, Theorem 9.9).

\begin{lemma} Let $n>2$ be an integer and let
$X$ denote a set with $n=|X|$.
$S_X$ acts $n$-transitively
on $X$ and $A_n$ acts $(n-2)$-transitively
on $\{1,2,...,n\}$.
\end{lemma}

\begin{lemma} (\cite{DM}, Corollary 1.4A)
\label{regular}
If $G$ acts transitively on $X$ then $G$ acts regularly
on $X$ if and only if $|G|=|X|$. 
\end{lemma}

The following general result might be regarded as our
main lemma. It's proof is omitted.

\begin{lemma} (\cite{DM}, Exercise 1.4.1)
If $G$ acts transitively on $X$, $x\in X$,
and $H$ is any subgroup of $G$ then
\[
G=G_xH=\{gh\ |\ g\in G_x,h\in H\}
\]
if any only if $H$ acts transitively on $X$. 
\end{lemma}

The following result gives precise conditions which may be
used to determine the answer to the question in the title
to this paper.

\begin{proposition} $S_m$ has a complement in $S_n$ if and 
only if there is an subgroup $H$ of $S_n$ such that
\begin{itemize}
\item
 $H$ acts $(n-m)$-transitively on $\{1,2,...,n\}$,
\item
$|H|=n(n-1)...(m+1)=n!/m!$.
\end{itemize}
\label{complement}
\end{proposition}

\begin{remark} The method of proof actually
proves something more general: one may replace $S_m$ in the 
statement of the above Proposition with any
stabilizer of $S_n$ acting on $\{1,2,...,n\}^{(n-m)}$. 
\end{remark}

{\bf proof}: Let $G=S_n$, which we regard as acting 
transitively on $X=\{1,2,...,n\}^{(n-m)}$.
Identify $S_m$ with the stabilizer $G_x$ of
$x=(m+1,...,n)\in \{1,2,...,n\}^{(n-m)}$. The above lemma says
that $G=G_xH$ for some $H$ if and only if $H$ acts 
transitively on $X$, i.e., 
acting $(n-m)$-transitively on $\{1,2,...,n\}$.
$H$ is a complement of $G_x$ if and only if
$H \cap G_x=1$, which is equivalent to saying that $H$ acts
regularly on $X$. Since $|X|=n!/m!$, the result follows from 
Lemma \ref{regular}. $\Box$

\begin{corollary} In general, $S_2=\langle (1,2)\rangle \subset S_n$
has a complement in $S_n$, namely $A_n$.
\end{corollary}

\begin{example} If $n=10$ and $m=6$ then the Theorem says that
$S_6$ has a complement in $S_{10}$ if and only if there is a subgroup 
$H$ of $S_{10}$ of order $5040$ which acts $4$-transitively
on $\{1,2,...,10\}$. The only subgroups of $S_{10}$
which act $4$-transitively on $\{1,2,...,10\}$ are $S_{10}$ itself and 
$A_{10}$. This follows from the classification of all 
$4$-transitive groups (see \cite{DM}, chapter 7, or
\cite{R}, page 289). This implies $S_6$ has no complement in
$S_{10}$. (However, $S_{10}$ does have a subgroup of
order $5040$ but it is not even transitive on
$\{1,2,...,10\}$.)

Similarly, $S_5$ has no complement in $S_{10}$. 

Note that $S_{10}$ has not one but two non-isomorphic subgroups,
$H_1$ and $H_2$,
of order $720=10!/7!$, each of which acts $3$-transitively
on $\{1,2,...,10\}$ (see \cite{DM}, chapter 7). These groups satisfy
$H_1\cong M_{10}$ and $H_2\cong PGL_2({{\Bbb F}}_9)$.
Thus $S_7$ has two non-isomorphic complements in $S_{10}$.
\end{example}

\begin{lemma} (\cite{DM}, pages 242-243) 
Let $q=p^k$, $p$ a prime and $k\geq 1$ an integer. 
There is a $2$-transitive group 
$G\cong AGL_1({{\Bbb F}}_q)$ acting on
$X={{\Bbb F}}_q$ of order $q(q-1)$.
Consequently, $S_{q-2}$ has a complement in $S_q$.
\label{2transitive}
\end{lemma}

\begin{lemma} (\cite{DM}, pages 242-243) 
Let $q=p^k$, $p$ a prime and $k\geq 1$ an integer. 
There is a $3$-transitive group 
$G\cong PGL_2({{\Bbb F}}_q)$ acting on
$X={{\Bbb F}}_q\cup \{\infty\}$ of order $(q+1)q(q-1)$.
Consequently, $S_{q-2}$ has a complement in $S_{q+1}$.
\label{3transitive}
\end{lemma}

The following Theorem is our main result.

\begin{theorem} 
\label{mainthrm}
\begin{itemize}

\item
If $n>2$ is not a prime power or
a prime power plus $1$ then the only $1<m<n$ for
which $S_m$ has a complement in $S_n$ are $m=2$
and $m=n-1$. 
\item
If $n>12$ is a prime power and not
a prime power plus $1$ then 
the only $1<m<n$ for which $S_m$ has a complement in 
$S_n$ are $m=2$, $m=n-2$ and $m=n-1$. 
\item
If $n>12$ is a prime power plus $1$ but not a prime power 
then the only $1<m<n$ for which $S_m$ has a complement in 
$S_n$ are $m=2$, $m=n-3$ and $m=n-1$. 
\item
If $n>12$ is both a prime power plus $1$ and a prime power 
then the only $1<m<n$ for which $S_m$ has a complement in 
$S_n$ are $m=2$, $m=n-3$, $m=n-2$ and $m=n-1$. 
\item
If $n\leq 12$ see the above table.
\end{itemize}
\end{theorem}

\begin{remark} Roughly speaking, this theorem says that if you 
pick an integer $n>2$ ``at random'' then, with
``probability one" (in the sense of Dirichlet
density), the only $S_m$'s
which have a complement in $S_n$ are $S_2$ and $S_{n-1}$.
\end{remark}

{\bf proof}: The proof below depends on the classification of 
$2$-transitive groups of degree $n$ (see \cite{DM}, 
\S\S 7.6-7.7). The following groups are $2$-transitive:
\begin{itemize}
\item
Some ``affine groups'' and some of their subgroups
are $2$-transitive. These groups
always act on a set of prime power degree.
\item
Projective special linear groups $PSL_d({{\Bbb F}}_q)$ 
acts $2$-transitively on projective $d-1$-space 
${{\Bbb P}}^{d-1}({{\Bbb F}}_q)$, a set of size
$(q^d-1)/(q-1)$. In general (see \cite{R},
chapter 8, page 223),
\[
|PSL_d({{\Bbb F}}_q)|={(q^d-1)(q^d-q)...(q^d-q^{d-1})
\over {gcd(d,q-1)\cdot (q-1)}}.
\]

In case $d=2$, $PGL_2({{\Bbb F}}_q)$ is sharply $3$-transitive
of degree $q+1$. If $q$ is even then 
$PGL_2({{\Bbb F}}_q)=PSL_2({{\Bbb F}}_q)$. If $q$ is odd
then $PSL_2({{\Bbb F}}_q)$ is a subgroup of index $2$ in
$PGL_2({{\Bbb F}}_q)$.

\item
The symplectic groups $Sp_{2m}({{\Bbb F}}_2)$, $m\geq 2$,
act $2$-transitively (in different ways) on a set of size 
$2^{m-1}(2^m+1)$ and a set of size $2^{m-1}(2^m-1)$
(see \cite{DM}, \S 7.7, page 245).
We have 
\[
|Sp_{2m}({{\Bbb F}}_2)|=2^{m^2}(2^2-1)(2^4-1)...(2^{2m } -1)
\]
(see \cite{R}, chapter 8, page 245).

\item
Special projective unitary groups $PSU_3({{\Bbb F}}_q)$.
These groups are of order
\[
|PSU_3({{\Bbb F}}_q)|=
{1\over gcd(3,q+1)}q^3(q^2-1)(q^3+1)
\]
(see \cite{R}, chapter 8, page 245)
and act $2$-transitively on a set of size $q^3+1$
(see \cite{DM}, \S 7.7, page 249).

\item
Suzuki groups $Sz(q)=\ ^2B_2({{\Bbb F}}_q)$,
where $q=2^{2m+1}$. This group is order
\[
|Sz(q)|=q^2(q-1)(q^2+1), \ \ q>2
\]
(see \cite{DM}, \S 7.7, page 250).
If $q=2$ then $Sz(2)\cong AGL_1({{\Bbb F}}_5)$.
This group acts $2$-transitively on the
$q^2+1$ points of a Steiner system of type $S(3,q+1,q^2+1)$.

\item
Ree groups $R(q)=\ ^2G_2({{\Bbb F}}_q)$,
where $q=3^{2m+1}$. This group is order
\[
|R(q)|=q^3(q-1)(q^3+1).
\]
(see \cite{DM}, \S 7.7, page 251).
This group acts $2$-transitively on the
$q^3+1$ points of a Steiner system of type $S(2,q+1,q^3+1)$.

\item 
The Mathieu groups (see \cite{R}, chapter 9):
\begin{itemize}
\item 
$M_{10}$ is $3$-transitive of degree $10$, 
$|M_{10}|=720$.
\item
$M_{11}$ is $4$-transitive of degree $11$, 
$|M_{11}|=7920$. Also, $M_{11}$ 
has a $3$-transitive action of degree $12$.
(See \cite{DM}, \S 7.7, page 252.)
\item
$M_{12}$ is $5$-transitive of degree $12$, 
$|M_{12}|=95040$.
\item
$M_{22}$ is $3$-transitive of degree $22$, 
$|M_{22}|=443520$.
\item
$M_{23}$ is $4$-transitive of degree $23$, 
$|M_{23}|=10200960$.
\item
$M_{24}$ is $5$-transitive of degree $24$,
$|M_{24}|=244823040$.  
\end{itemize}

\item
$PSL_2({{\Bbb F}}_{11})$ has a $2$-transitive action of degree
$11$ (recall $|PSL_2({{\Bbb F}}_{11})|=120\cdot 110/20 = 660$).
(See \cite{DM}, \S 7.7, page 252.)

\item
$A_7$ has a $2$-transitive action of degree $15$
(see \cite{DM}, \S 7.7, page 252).

\item
Higman-Sims group $HS$ of order 
$44352000$ has a $2$-transitive action on a set of size $176$
(see \cite{DM}, \S 7.7, pages 252-253).

\item
The Conway group $Co_3$ of order $495766656000$
has a $2$-transitive action of degree $276$
(see \cite{DM}, \S 7.7, page 253).
\end{itemize}

There are no other $2$-transitive groups. 

To see if $S_m$ has a complement in $S_n$ not already listed
(i.e., to try to find a counterexample to the statement
of the Theorem above), we must find an $(n-m)$-transitive
group of degree $n$ and order $n!/m!$. Because of
Lemmas \ref{2transitive} and \ref{3transitive},
we may assume
$m\in \{2,3,4,5\}$, $n>12$, $(n,m)\not=({\rm a\ prime\ power},n-2)$,
and $(n,m)\not=({\rm a\ prime\ power\ plus}\ 1,n-3)$.

The argument shall be case-by-case:

Case 1: The projective special linear groups 
$PSL_d({{\Bbb F}}_q)$.
We may assume $d>2$ since otherwise $n$ will be
a prime power plus $1$.
In this case $n=(q^d-1)/(q-1)$ and $m=n-2$. In order
for $H=PSL_d({{\Bbb F}}_q)$ to be a complement
of $S_{n-2}$ we must have
\[
|H|=|PSL_d({{\Bbb F}}_q)|={(q^d-1)(q^d-q)...(q^d-q^{d-1})
\over {gcd(d,q-1)\cdot (q-1)}}=n(n-1),
\]
by Proposition \ref{complement}. In other words, $q$ must satisfy
\[
{(q^d-q)...(q^d-q^{d-1})
\over {gcd(d,q-1)}}=n-1=(q^d-q)/(q-1),
\]
which is impossible if $d>2$.

Case 2: The symplectic groups $Sp_{2m}({{\Bbb F}}_2)$, $m\geq 2$.
In this case $n=2^{r-1}(2^r+1)$ or $n=2^{r-1}(2^r-1)$ 
and $m=n-2$. In order
for $H=Sp_{2r}({{\Bbb F}}_2)$ to be a complement
of $S_{n-2}$ we must have
\[
|H|=|Sp_{2r}({{\Bbb F}}_2)|=2^{r^2}(2^2-1)(2^4-1)...(2^{2r } -1)
=n(n-1).
\]
In other words, $r$ must satisfy either
\[
2^{r^2}(2^2-1)(2^4-1)...(2^{2r } -1)
=2^{r-1}(2^r+1)(2^{r-1}(2^r+1)-1)
\]
or
\[
2^{r^2}(2^2-1)(2^4-1)...(2^{2r } -1)
=2^{r-1}(2^r-1)(2^{r-1}(2^r-1)-1).
\]
Both of these equalities are impossible since
for $r>1$ the left-hand side is always larger than
the right-hand side.

Case 3: The special projective unitary groups 
$PSU_3({{\Bbb F}}_q)$
In this case $n=q^3+1$ and $m=n-2$. In order
for $H=PSU_3({{\Bbb F}}_q)$ to be a complement
of $S_{n-2}$ we must have
\[
|H|=|PSU_3({{\Bbb F}}_q)|=
{1\over gcd(3,q+1)}q^3(q^2-1)(q^3+1)
=n(n-1)=(q^3+1)q^3.
\]
This is impossible unless $gcd(3,q+1)=q^2-1$, i.e., unless
$q=2$. But if $q=2$ then the degree is $n=9$ and we have
already covered this case in the table above.

Case 4: The Suzuki groups $Sz(q)=\ ^2B_2({{\Bbb F}}_q)$,
where $q=2^{2m+1}$. In this case $n=q^2+1$ and $m=n-2$. In order
for $H=Sz(q)$ to be a complement of $S_{n-2}$ we must have
\[
|H|=|Sz(q)|=q^2(q-1)(q^2+1)
=n(n-1)=(q^2+1)q^2.
\]
This is impossible.

Case 5: The Ree groups $R(q)=\ ^2G_2({{\Bbb F}}_q)$,
where $q=3^{2m+1}$. In this case $n=q^3+1$ and $m=n-2$. In order
for $H=R(q)$ to be a complement of $S_{n-2}$ we must have
\[
|H|=|R(q)|=q^3(q-1)(q^3+1)
=n(n-1)=(q^3+1)q^3.
\]
This is impossible.

Case 6: The Mathieu groups. Since we are assuming $n>12$, we need 
only consider
\begin{itemize}
\item
$M_{22}$ is $3$-transitive of degree $22$, 
$|M_{22}|=443520$.
\item
$M_{23}$ is $4$-transitive of degree $23$, 
$|M_{23}|=10200960$.
\item
$M_{24}$ is $5$-transitive of degree $24$,
$|M_{24}|=244823040$.  
\end{itemize}
None of these satisfy $|H|=n!/m!$, for $m=n-3,n-4,n-5$, 
respectively.

Case 7: $A_7$ has a $2$-transitive action of degree $15$.
In this case, $n=15$ and $m=n-2$. But
$|A_7|=2520\not= 15\cdot 14=210$, so $A_7$ is not 
a complement of $S_{13}$ in $S_{15}$.

Case 8: The Higman-Sims group $HS$. In this case $n=176$ and
$m=n-2$. But
$|HS|=44352000\not= 176\cdot 175= 30800$, so $HS$ is not 
a complement of $S_{174}$ in $S_{176}$.

Case 9: The Conway group $Co_3$. In this case, 
$n=276$ and $m=n-2$. But
$|Co_3|=495766656000\not= 276\cdot 275=  75900$, so $Co_3$ is not 
a complement of $S_{274}$ in $S_{276}$. 

This exhausts the $2$-transitive groups, so the proof is complete.

$\Box$

\section{Final remarks}

From the table in \S 2, we see that $S_8$ does not have a complement in 
$S_{12}$. However, $PGL_2({{\Bbb F}}_{11}), C_9, S_8$ are complementary
in $S_{12}$ (in the sense of Definition \ref{def:k-compl}),
\[
S_{12}=PGL_2({{\Bbb F}}_{11})\cdot C_9\cdot S_8.
\]
This yields one method of storing permutations in $S_{12}$
using those of $S_8$.

More generally, given any $m$ and any $n$ with $m<n$, Theorem \ref{mainthrm}
and the table in \S 2 may be used to find the smallest $k$ such
$G_1$, $G_2$, ..., $G_k=S_m$ are complements (in the sense of Definition
\ref{def:k-compl}) in $S_n$.

{\it Acknowledgements}: I'd like to thank Herbert Kociemba for 
several interesting and inspiring communications. This paper
used GAP \cite{GAP} to verify some of the results stated in
the table in \S 2.


\begin{thebibliography}{99}

\bibitem[C]{C}
Peter Cameron, {\bf Permutation Groups}, 
Cambridge University Press, 1999.

\bibitem[CHM]{CHM}
John Conway, Alexander Hulpke, and John McKay.
{\it On transitive permutation groups},
LMS J. Comput. Math., 1 (1998), 1-8.

\bibitem[DM]{DM} John Dixon, Brian Mortimer, 
{\bf Permutation groups}, Springer-Verlag, 1996.

\bibitem[GAP]{GAP} The GAP Group, 
{\it GAP -- Groups, Algorithms, and Programming}, 
Version 4.4.9; 2006. (http://www.gap-system.org).

\bibitem[K]{K} Herbert Kociemba, private communications, 
Summer 1997. His cube explorer program is at
http://kociemba.org/.

\bibitem[R]{R} Joseph Rotman, 
{\bf An introduction to the theory of groups}, 
$4^{th}$ edition, Springer-Verlag, 1995.

\end{thebibliography}

\end{document}